\bookchapter{Occupied Atoms from the Complete Nuclear Source}

\subsection{From a neutral count to an electronic state}

A neutral ledger specifies an electron inventory, $N_e=Z$. To construct an atom one must also specify the electronic space, its kinetic action, the field to which it couples, its antisymmetry, and the energy of the separated constituents. The completed Volume II work supplies these objects on the original helium and carbon sources. The transition is substantial: the electron entry becomes an occupied state with orbital coefficients, spin labels, radial moments, a complete Hamiltonian and an explicit ionization reference. Nuclear phase multiplicity remains part of the object throughout.

The helium inputs contain $2p2n$ and all $4096$ original phase histories per cover. The carbon inputs contain $6p6n$ and all $68{,}719{,}476{,}736=2^{36}$ original histories. Their intrinsic energies and signed receiver $U$ are retained. At the ground-state level these sources have $49$ helium conditional receiver sectors and $361$ ordered carbon sectors. A shared electron source acts on each of these original nuclear states before the nuclear minimum or Gibbs average is taken. The construction consequently permits electronic occupation to change which nuclear sector is energetically selected. \samref{electronic_source}\samref{radial_occupancy}

The initial electronic Hilbert space is
\[
\begin{gathered}
\mathcal H_1=L^2(S^2)\otimes\mathbb C^2,
 \\
 \mathcal H_{N_e}=\bigwedge^{N_e}\mathcal H_1.
\end{gathered}
\]
The two components are declared spin labels. Antisymmetrization determines the available occupations. Helium retains every charge sector $N_e=0,1,2$; carbon retains $N_e=0,\ldots,6$. On a sphere of radius $R$ the one-electron generator is
\[
\begin{gathered}
h_U=\frac{\tau}{R^2}(-\Delta_{S^2})-\frac gR U\cdot n,
 \\
 H(q)=E_{\rm nuc}(q)+d\Gamma(h_{U(q)}).
\end{gathered}
\]
Here $d\Gamma$ is the fermionic sum of the one-body action and $n$ is the position direction. The coefficient $\tau$ controls kinetic action. The same actual $U(q)$ couples to every electron. Replacing that common nuclear variable by independent electron averages would define another source.

This electronic law, its radial continuations and the field incidence developed below are Codex mathematical contributions, consuming the nuclear constructions directed by Sean Brady. The completed source calculations use the project's classification: \emph{The test result suggests strong contact with the concept.} The physical charge, energy and length assignment is carried by its own common-source recipient; it is not supplied by the neutral count.

\subsection{The full angular trace and its preparation payment}

The angular construction keeps the kinetic operator inside the quantum exponential. If $p_0$ is the complete original nuclear Gibbs preparation, its canonical electronic trace is
\[
\begin{gathered}
Z_{N_e}=\mathbb E_{p_0}\operatorname{Tr}_{\mathcal H_{N_e}}
           e^{-\beta d\Gamma(h_U)},\\
 \Delta F_{N_e}=-\frac1\beta\log\frac{Z_{N_e}}{Z^{\rm free}_{N_e}}.
\end{gathered}
\]
The free reference has the same radius, temperature and electron number and includes the complete angular spectrum. Its rotational symmetry makes the initial interaction work zero. The joint relative entropy then supplies the relaxation payment
\[
 D(\rho_0\Vert\rho)=-\beta\Delta F_{N_e}.
\]
For nonzero coupling and nonempty occupation the complete-source free-energy change is strictly negative; the empty sector gives exactly zero. This connects a prepared atomic response to an explicit thermodynamic ledger. \samref{electronic_source}

Rotational covariance reduces each conditional spectrum to a function of $|U|^2$, while the original nuclear weights stay attached to that receiver. The carbon source has $51$ receiver-squared values. Across the declared parameter panel, exact sharing yields $240$ dimensionless spectral packets. At angular prefix $L=32$, each packet retains $1089$ spatial modes and $2178$ spin orbitals. The number of possible $N_e$-electron occupations is a binomial count in that orbital inventory; the calculation does not form a matrix of that binomial dimension.

The angular complement is retained by an operator bound. With $P=P_L$, $Q=1-P$, $a=g|U|/R$ and kinetic operator $K$, the omitted link satisfies
\[
\begin{gathered}
b^2\leq\frac{a^2(L+1)^2}{(2L+1)(2L+3)},
 \\
 h_U\geq(P h_U P-\delta)\oplus
          (QKQ-a-b^2/\delta),\quad\delta>0.
\end{gathered}
\]
Finite-sector traces give one partition endpoint; this inequality and the infinite free tail give the other. The completed panel contains $1620$ electron-sector accounts. All $264$ populated $L=32$ accounts have negative full-source $\Delta F_{N_e}$ bounds, hence positive relaxation payments. All $60$ empty accounts remain zero.

The shared source also gives a concrete departure from independent-electron addition. For the connected readout $\Delta F_2-2\Delta F_1$, $37$ cohorts have positive intervals, five have negative intervals and eighteen contain zero. In the $\beta=R=g=1$ controls, helium cover 3 has positive connected response, whereas cover 4 has negative response. For cover 3, the displayed one- and two-electron intervals are approximately $[-0.233294,-0.232789]$ and $[-0.433514,-0.432425]$; for cover 4 they are $[-0.551141,-0.549727]$ and $[-1.129610,-1.126065]$. At this stage the response combines shared nuclear averaging and Pauli occupation. The direct pair field enters subsequently.

\subsection{Radial binding and the occupied orbital family}

The three-dimensional completion uses
\[
 h_U=\tau(-\Delta_{\mathbb R^3})-g\frac{U\cdot n}{r}.
\]
A free electron has continuum ground infimum zero. The separated reference is therefore the original nuclear ground plus the specified number of free electrons. This ground account and the finite-temperature sphere account retain their respective preparations. Completing the full gradient square gives the whole-space lower bound $-g^2|U|^2/(4\tau)$ per electron. A regular two-state $s/p$ trial has kinetic matrix $(\tau/a^2)I$ and field cross term $-g|U|/(3a)$. Its lower Ritz level selects
\[
\begin{gathered}
a=\frac{6\tau}{g|U|},\\
 E_{\rm trial}=-\frac{g^2|U|^2}{36\tau}.
\end{gathered}
\]
Two opposite spin labels may occupy that spatial level. The exact two-electron exterior action retains all six antisymmetric states of the initial trial space. Minima over every original nuclear sector yield twenty positive populated one- and two-electron binding accounts. \samref{electronic_source}

For general occupation, rotate $U$ to the polar axis and put $A=g|U|$, $r=(\tau/A)\rho$. The operator becomes
\[
 h_U=\frac{A^2}{\tau}\left[-\Delta_\rho-\frac{\cos\theta}{\rho}\right].
\]
One universal orbital calculation consequently serves every original receiver, with $U=0$ retaining the free continuum. Regular $r^l e^{-r/a}Y_{lm}$ multiplets retain every signed $m$ and both spins. Four radial functions per angular mode, with relative lengths $1/2,1,2,4$, form physical Gram pencils through $L=16$. The largest individual pencil block is $68$, and the full $L=16$ family has $2312$ spin orbitals. \samref{radial_occupancy}

Six actual occupied spatial shapes are returned with normalized coefficients and radial moments. Carbon's best six-electron trial occupies two distinct $m=0$ levels and the $m=\pm1$ multiplet. For fixed coefficients let $T_{N_e}$ and $C_{N_e}$ be the complete occupied kinetic and field expectations. Common dilation gives
\[
\begin{gathered}
E(a)=\frac{\tau T_{N_e}}{a^2}-\frac{A C_{N_e}}a,
 \\
 a=\frac{2\tau T_{N_e}}{A C_{N_e}},
 \\
 E_{\rm opt}=-\frac{A^2 C_{N_e}^2}{4\tau T_{N_e}}.
\end{gathered}
\]
All six occupied coefficients improve under this dilation, and their full virial residual $2K+V$ contains zero. The dimensionless coefficients are approximately
\[
\begin{array}{c|rrr}
 N_e&1&2&3\\\hline
 E_{\rm opt}\tau/A^2&-0.04720782&-0.09441564&-0.11456431\\[3pt]
 N_e&4&5&6\\\hline
 E_{\rm opt}\tau/A^2&-0.13500524&-0.14862815&-0.16225149
\end{array}
\]

All forty-four populated whole-source binding lower bounds are positive. At $\tau=g=1$, the six-electron carbon binding intervals are approximately $[5.84105,66]$, $[10.38410,104]$ and $[23.36421,216]$ for the 333, 343 and 444 sources. These source-energy intervals already include the original nuclear minimization. They show explicitly how a complete electronic state can sit on the same retained nuclear history rather than being appended as an integer label.

\subsection{Electronic occupation feeds back into the nuclear minimum}

The universal full-space electronic ground coefficient $e_{N_e}$ obeys
\[
 -N_e/4\leq e_{N_e}\leq c_{N_e}<0.
\]
The exact coupled nuclear selection is the lower envelope
\[
\begin{gathered}
\min_q\{E_{\rm nuc}(q)-tq\},\\
 t=(-e_{N_e})g^2/\tau,\\
 q=|U|^2.
\end{gathered}
\]
Every original nuclear sector participates in this envelope. The completed calculation obtains eleven response phases and six exact transition values. Helium cover 3 changes from $q=4$ to $36$ at $t=43/8$; cover 4 changes from $16$ to $36$ at $129/5$. Carbon 333 changes from $36$ to $324$ at $35/24$; carbon 343 follows $64\to256\to324$ at $35/24$ and $123/17$; carbon 444 changes from $144$ to $324$ at $124/15$. \samref{radial_occupancy}

Combining the full electronic coefficient bounds with this exact nuclear envelope supplies forty-four nonempty coupling ranges that guarantee a nuclear receiver phase for the whole-space electron ground. The intermediate carbon 343 envelope exists exactly; its six intermediate coupling guarantees are unresolved with those electronic bounds. This is a precise remaining assignment within an otherwise completed response calculation. The envelope retains the nuclear energy price of changing receiver, so a larger $|U|$ is not selected freely.

This feature provides an important construction principle for the rest of the volume. The source, recipient and occupied state are solved together through a shared energy. A signed receiver is neither a fixed empirical charge nor a coefficient that can be transferred independently of its nuclear energy. Subsequent scalar and polarization fields preserve that distinction while adding explicit interactions and finite core self energy.

\subsection{One scalar and three polarization channels}

The interacting atom uses the Green normalization $G(x,y)=1/(2\pi|x-y|)$ and a positive stored constrained-displacement energy. A scalar channel assigns proton $+1$, neutron $0$ and electron $-1$. Three polarization channels retain the original vector $U$ and the electron direction $n(x)$. The incidence is a declared source construction. For a uniform nuclear shell of radius $r_c$, the resulting terms are
\begin{align*}
 S_{\rm core}(q)&=\frac{g_c^2Z^2/v_c+g_p^2q/v_p}{2r_c},\\
 V_{\rm core}(x)&=-\frac{g_cZ+g_pU\cdot n(x)}{\max(r,r_c)},\\
 W(x,y)&=\frac{v_c+v_p n(x)\cdot n(y)}{|x-y|}.
\end{align*}
The coefficients derive from the source amplitudes and field stiffness: $g=\lambda_{\rm core}\lambda_e/(2\pi\kappa)$ and $v=\lambda_e^2/(2\pi\kappa)$. Equal scalar proton/electron magnitudes give $g_c=v_c$. For $v_c\geq v_p\geq0$ the pair interaction is pointwise nonnegative. \samref{pair_derivation}\samref{pair_source}

Core self energy stays inside both nuclear minima. The regularized electron self, $N_e(v_c+v_p)/(2r_e)$, is common at fixed inventory; it cancels between coupled and separated preparations before taking the point-electron limit. A $q$-dependent core self cannot be cancelled between different nuclear states. This bookkeeping determines the full binding reference:
\[
\begin{gathered}
E_{\rm sep}=\min_q[E_{\rm nuc}(q)+S_{\rm core}(q)],\\
 E_{\rm atom}=\min_q[E_{\rm nuc}(q)+S_{\rm core}(q)+E_{e,N_e}(q)].
\end{gathered}
\]
Binding is $E_{\rm sep}-E_{\rm atom}$. Both use the same field recipient and intrinsic source.

The retained radial orbitals generate every product multipole through rank $32$. The polarization contraction also carries adjacent radial orders through $33$. Exact Gaunt coefficients and complete two-region radial integrals retain off-diagonal orbital products and signed magnetic terms. These products define direct and exchange entries on the actual physical Gram. The calculation therefore includes the changing pair field of the occupied state, with its own angular structure.

The complete retained fermion spaces use one, two and four spatial orbitals with both spins. All configurations through the actual helium and carbon inventories are retained. Grouping by particle number, spin-up inventory and total magnetic number preserves every configuration. The resulting $59{,}052$ certified configuration spectra include kinetic action, finite-core attraction and full scalar/polarization pair terms. Their largest block has twelve states. The original all-angular and all-radial complement is covered separately by the whole-space lower account.

\subsection{Pauli bounds, ion thresholds and complete filling}

The antisymmetric lower account fills distinct Coulomb spin states. For $N_e\leq6$ its occupation weight is
\[
 w_{N_e}=\min(N_e,2)+\frac{\max(N_e-2,0)}4.
\]
For six electrons this gives $3$, improving the distinguishable occupation weight $6$. With $A(q)=g_cZ+g_p\sqrt q$, the full-space electronic lower bound is
\[
 E_{e,N_e}(q)\geq
 \max\left\{-\frac{w_{N_e}A(q)^2}{4\tau},
                  -\frac{N_eA(q)}{r_c}\right\}.
\]
The first entry is the complete Pauli spectral comparison; the second uses the finite-core potential floor. Positive pair energy may be discarded for this lower account, whereas the upper witness retains its full pair and exchange action. \samref{atomic_derivation}\samref{pair_source}

The completed interacting panel has positive total binding lower bounds in all forty-four populated accounts. Its initial upper portfolio sometimes selected a localized ion plus free electrons. All ion/free thresholds were retained, yielding ten positive numerical ionization lower bounds and thirty-four intervals containing zero. The next construction supplies fully occupied witnesses in every populated account and an analytic retention result on the same Hamiltonian. \samref{outer_filling}

A spherical outer orbital has support $[R,2R]$ and radial polynomial $s(1-s)P_4(s)$, $s=r/R-1$. The four coefficients enforce orthogonality to all four actual inner $l=0$ radial primitives. Angular orthogonality handles the other modes; later shells have disjoint support. The correlated inner configuration-interaction state remains intact. Exchange subtracts a positive quadratic form from the direct upper account. For electron $j$, the complete exterior bound has the form
\[
 \Delta E_j\leq\frac{\tau K+v_pM_1J_2/3}{R^2}
                  -\frac{[g_cZ-v_c(j-1)]J_1}{R}.
\]
Positive residual scalar charge makes the $1/R$ attraction dominate. Thus the whole-space ground decreases strictly on adding electron $j$ whenever $g_cZ>v_c(j-1)$, with the stated positive parameters and $v_c\geq v_p$. At $g_c=v_c$ every $j\leq Z$ qualifies. Every neutral helium and carbon sector therefore retains all its electrons, including every original nuclear phase.

The added shells replace nine earlier ion/free witnesses and improve ten upper energies. For carbon 333 at $g_p=v_p=1/10$, the five- and six-electron improvements are approximately $0.00244396$ and $0.00389937$ source units; the latter uses outer radii $64$ and $256$. The analytic filling result covers all forty-four populated accounts. The quantitative numerical gap ledger remains ten positive intervals and thirty-four overlaps. These two completed outputs answer different questions: existence of strict retention and the current enclosed size of that retention energy.

\subsection{A worked interacting binding account}

Consider the weak-polarization panel $\tau=r_c=g_c=v_c=1$, $g_p=v_p=1/10$. The original interacting helium cover-3 and cover-4 neutral accounts both return a total binding interval approximately $[0.436748933182,2.22]$. Their equality in this interacting panel is a result of the complete source minima. It is obtained after finite core self and the occupied electronic action have been added, rather than by setting the two receivers equal. In the stronger polarization panel $g_p=v_p=1$, their retained total-binding interval is approximately $[0.368497146536,6]$. The same upper-reference rule is applied in each case. \samref{pair_source}

For carbon 444 with weak polarization, the interacting neutral interval is approximately $[6.0992574809,38.88]$. At that stage its best upper witness localized five of the six electrons. The complete account still contained six electrons: one occupied the separated free-electron continuum. The compact exterior-shell construction subsequently replaced this witness by a normalized six-electron determinant joined to the correlated inner state, lowering its upper energy by about $0.00126403124427$. This is a useful example of how the ledger records a physical improvement. The total reference is held fixed, an explicit additional localized state is supplied, and its complete extra kinetic, direct and exchange cost is calculated. \samref{outer_filling}

The electron count alone would not reveal that difference between the two witnesses. The actual state specifies it through its support and occupations. The finished neutral atom has a fully occupied variational witness, and the analytic residual-charge construction establishes strict whole-space retention. The numerical ionization interval continues to report the available quantitative gap enclosure. A future tighter interval can be compared with the same reference without revising the meaning of the earlier result.

A related distinction applies to an orbital radius. The radial dilation $a$, a mean radius, a compact shell's support radius and the nuclear field-shell radius $r_c$ are separately stored quantities. The original radial trial has mean radius $2a$ and mean-square radius $21a^2/4$; an outer shell occupies $[R,2R]$. Neither is automatically a measured atomic radius. The completed constructions retain enough information to define a radius-sensitive receiver later, with its actual source normalization and energy cost.

The execution record follows these constructions at their stated scope. The first complete electronic source used $422$ native scientific returns and $105{,}924$ independent assertions. The interacting pair and configuration calculation used $2237$ native returns and $16{,}730$ independent assertions; the full filling continuation added $148$ native returns and $3026$ checks without repeating the earlier pair spectra. These counts identify completed operands and their independent reconstructions. They are not additional physical assumptions. Full copied-session recovery retained the published values, including the corrected pair calculation and the later outward endpoint correction.

\subsection{Radial freedom in helium, carbon and oxygen}

The later He/C/O refinement brings the same construction into a twenty-four-spin-orbital preparation family. Ten inherited spatial orbitals remain, including the original $264$-component source. Two compact/diffuse $m=0$ seeds, with radial lengths $1/2$ and $2$, are projected off the existing orbitals and each other in the full physical Gram. Each new row has sixty-eight components; the complete twelve-orbital preparation has $406$. The Gram enclosure contains the identity. \samref{radial_pod}\samref{radial_correction}

Fifty-four matched controls cover three species, three dilations, three radii and two inventory references. Every charge $N_e=0,\ldots,Z$ and every original neutral phase remains attached to its readout. The twenty-spin family is a literal subset. Neutral oxygen now has $735{,}471$ Slater configurations per control, and the entire expanded comparison makes $26{,}315{,}604$ all-charge occupation comparisons. These are complete occupations of the declared family.

Forty-five of fifty-four neutral upper witnesses improve and lie below the entire corrected preceding band. All fifty-four neutral lower readouts improve. In the inherited source-to-reference normalized eV readout, best attachment lower margins change from $21.85869$ to $54.90420$ for helium, $445.48578$ to $876.34827$ for carbon and $824.36874$ to $1651.35979$ for oxygen. Each best witness uses core radius $1/65536$; oxygen's best dilation changes from $1/2$ to $1$. These are prepared-family attachment margins on the specified free-inventory reference. Full configuration mixing, independently relaxed ion states and the common physical-unit assignment remain their named recipients.

The whole massive action, every same-spin exchange term, full collective recoil and momentum outside the orbital span are retained. Exact angular-sector assembly also reproduces the complete Gram and second/fourth moments through fifty-five small angular sectors. All $432$ difference enclosures contain zero against the primary route. The source's physical content survives this computational rearrangement: the same coefficient map reconstructs all $406$ components.

The preceding one-body upper-endpoint selection was corrected without changing native primitives. Rebuilding $252$ controls and $4788$ endpoints raises an upper by at most $8.281\times10^{-20}$ source units; lower endpoints and winning occupations are unchanged. The retained correction is part of the finished result. Together, the angular trace, occupied radial orbitals, scalar/polarization pair source, charge-filling construction and radial improvement supply an actual atomic chapter with explicit states, energies and references. Their authorized classification is strong contact; their remaining physical assignments are the common energy/action/ruler map and the specifically identified ground and ion recipients.
